\section{Matrix identities, isometric frames and nonnegative matrices}

\begin{definition}[Embedding replacing an occupied frame]
    \label{def:isometric_frame_embedding}
    \lean{Matrix.isometricFrameEmbedding, Matrix.finIsometricFrameEmbedding,
        Matrix.unitaryDirectSum, Matrix.unitaryDirectSumHom}
    \leanok
    Let $F:S\to H$ and $J:S\to K$ be linear isometries between finite-dimensional
    complex Hilbert spaces, and let $E:H\to K\oplus H$ be given by
    \begin{align}
        \label{eq:isometric_frame_embedding}
        Ev=(JF^\dagger v,(I-FF^\dagger)v).
    \end{align}
    The first summand carries the prescribed target frame, while the second
    retains the orthogonal complement of the range of $F$. For unitaries $V$ on
    $K$ and $U$ on $H$, write $V\oplus U$ for their direct sum on $K\oplus H$.
    This places two isometric frames in a common space, as
    in~\cite[Section~II.F.2]{Schuch2011Phases}.
\end{definition}

\begin{lemma}[Isometry and covariance of the frame embedding]
    \label{thm:isometric_frame_embedding}
    \lean{Matrix.isometricFrameEmbedding_isometry,
        Matrix.isometricFrameEmbedding_mul_frame,
        Matrix.conjTranspose_intertwiner_of_unitary,
        Matrix.isometricFrameEmbedding_intertwiner,
        Matrix.finIsometricFrameEmbedding_isometry,
        Matrix.finIsometricFrameEmbedding_intertwiner,
        Matrix.finIsometricFrameEmbedding_mul_frame}
    \leanok
    \uses{def:isometric_frame_embedding}
    The map $E$ in~(\ref{eq:isometric_frame_embedding}) satisfies $E^\dagger E=I$
    and $EF=(J,0)$. If unitary actions obey $U_gF=FS_g$ and $V_gJ=JS_g$,
    then $(V_g\oplus U_g)E=EU_g$. The direct sum preserves multiplication,
    so two unitary representations give a unitary representation on $K\oplus H$.
\end{lemma}

\begin{proof}\leanok
    \uses{def:isometric_frame_embedding}
    With $P=FF^\dagger$, idempotence and self-adjointness $P^2=P=P^\dagger$
    give $E^\dagger E=FJ^\dagger JF^\dagger+(I-P)^2=P+I-P=I$ and $EF=(J,0)$.
    Taking adjoints of $U_gF=FS_g$ and multiplying by the unitary actions
    gives $F^\dagger U_g=S_gF^\dagger$. Thus $U_gP=PU_g$ and
    $V_gJF^\dagger=JF^\dagger U_g$, establishing the two components of the
    intertwining identity.
\end{proof}

\begin{lemma}[Inverse of a unitary invertible matrix]
    \label{lem:unitary_general_linear_inverse}
    \lean{Matrix.coe_gl_inv_eq_conjTranspose_of_mem_unitaryGroup}
    \leanok
    The inverse of an invertible complex matrix that is unitary is its
    conjugate transpose.
\end{lemma}

\begin{proof}\leanok
    The conjugate transpose is a left inverse, and a left inverse of an
    invertible element is its inverse.
\end{proof}

\begin{lemma}[Coordinate inclusions]
    \label{lem:coordinate_inclusion}
    \lean{Matrix.coordinateInclusion,
        Matrix.coordinateInclusion_isometry,
        Matrix.conjTranspose_coordinateInclusion_mul,
        Matrix.mul_coordinateInclusion,
        Matrix.coordinateInclusion_compression,
        Matrix.coordinateInclusion_projection_mul_of_rowSupport,
        Matrix.coordinateInclusion_intertwine_of_rowSupport}
    \leanok
    An injection $e:\{1,\ldots,D\}\to\{1,\ldots,K\}$ defines an
    isometry $J_e$ by $J_e\ket a=\ket{e(a)}$. Its adjoint selects the
    included coordinates, so $J_e^\dagger M J_e$ is the corresponding
    submatrix. If the rows of $M$ are supported on these coordinates,
    then $J_eJ_e^\dagger M=M$ and $MJ_e=J_e(J_e^\dagger M J_e)$.
\end{lemma}

\begin{proof}\leanok
    The columns of $J_e$ are distinct coordinate vectors. Multiplication
    by $J_e$ and $J_e^\dagger$ selects columns and rows.
\end{proof}

\begin{lemma}[Anticommuting invertible operators force even dimension]
    \label{lem:even_finrank_of_anticommute}
    \lean{Module.End.even_finrank_of_anticommute,
        Module.End.isUnit_restrict_of_isUnit,
        Module.End.even_finrank_eigenspace_of_anticommute,
        Matrix.even_finrank_eigenspace_of_anticommute}
    \leanok
    Let $f,g$ be invertible operators on a finite-dimensional vector
    space $V$ over a field of characteristic zero with $fg=-gf$. Then
    $\dim V$ is even. Consequently, if a matrix $\Lambda$ commutes with two
    anticommuting invertible matrices $X$ and $Y$, every eigenspace of
    $\Lambda$ has even dimension.
\end{lemma}

\begin{proof}\leanok
    Taking determinants in $fg=-gf$ gives
    $\det f\det g=(-1)^{\dim V}\det g\det f$ with $\det f\det g\neq0$, so
    $(-1)^{\dim V}=1$. For the second statement, $X$ and $Y$ map each eigenspace
    of $\Lambda$ into itself, and their restrictions are invertible and
    anticommute.
\end{proof}

\begin{definition}[Extension of a matrix along an isometry]
    \label{def:spt_supported_matrix_extension}
    \lean{Matrix.isometryExtendByIdentity}
    \leanok
    For an isometry $K:\C^D\to\C^k$ and a matrix $X\in\MN{D}$, set
    $\widetilde X=KXK^\dagger+(I-KK^\dagger)$.
\end{definition}

\begin{lemma}[Comparison of supported frames]
    \label{thm:spt_supported_frame_comparison}
    \lean{Matrix.isometryExtendByIdentity_mul_isometry,
        Matrix.isometry_conjTranspose_mul_isometryExtendByIdentity,
        Matrix.isometry_reconstruct_of_supportProj,
        Matrix.commute_isometryExtendByIdentity_of_commute_compression,
        Matrix.isometry_absorption_of_projector_absorption,
        Matrix.isIsometry_conjTranspose_mul_of_absorption,
        Matrix.isometry_double_compression_eq_of_absorption,
        Matrix.commute_compression_of_commute_isometryExtendByIdentity,
        Matrix.supportFrame_dimension_pos_of_trace_one}
    \leanok
    \uses{def:spt_supported_matrix_extension}
    Let $K:\C^D\to\C^k$ be isometric. Then
    $\widetilde X K=KX$ and $K^\dagger\widetilde X=XK^\dagger$.
    If $\sigma\geq0$ has support $KK^\dagger$, then
    $\sigma=K(K^\dagger\sigma K)K^\dagger$; commutation of $X$ with
    $K^\dagger\sigma K$ implies commutation of $\widetilde X$ with $\sigma$.
    If $J:\C^r\to\C^k$ is isometric and $KK^\dagger J=J$, then
    $K^\dagger J$ is isometric and
    $(K^\dagger J)^\dagger(K^\dagger M K)(K^\dagger J)=J^\dagger M J$.
    The absorption $KK^\dagger JJ^\dagger=JJ^\dagger$ already implies
    $KK^\dagger J=J$.
    An operator $P$ commuting with $\widetilde X$ has compression
    $K^\dagger P K$ commuting with $X$.
    Finally, $\operatorname{tr}\sigma=1$ forces $D>0$.
\end{lemma}

\begin{proof}\leanok
    Multiply the defining matrix extension by $K$ and $K^\dagger$ and
    use $K^\dagger K=I$. The support projection absorbs $\sigma$ on
    both sides, giving the reconstruction identity. The same identities
    prove the commutation and nested-compression assertions.
    Multiplication of $KK^\dagger JJ^\dagger=JJ^\dagger$ by $J$ gives
    frame absorption. If $D=0$, the support is zero and hence
    $\sigma=0$, contradicting its trace.
\end{proof}

\begin{lemma}[Matrices of orthogonal projections]
    \label{lem:star_projection_matrices}
    \lean{Submodule.starProjectionMatrix,
        Submodule.toEuclideanLin_starProjectionMatrix,
        Submodule.isStarProjection_starProjectionMatrix,
        Matrix.isStarProjection_reindex,
        Matrix.isStarProjection_toEuclideanCLM_iff,
        Matrix.isSymmetricProjection_toEuclideanLin_iff}
    \leanok
    The matrix, in the standard basis, of the orthogonal projection onto a
    subspace of $\C^n$ is self-adjoint and idempotent, and a matrix is
    self-adjoint and idempotent exactly when the operator it defines is an
    orthogonal projection. Relabelling both
    coordinates of a self-adjoint idempotent matrix along one bijection gives
    a self-adjoint idempotent matrix.
\end{lemma}

\begin{proof}\leanok
    The passage from operators to matrices is a star-algebra isomorphism, and
    so is relabelling the coordinates; both preserve self-adjointness and
    idempotence.
\end{proof}

\begin{lemma}[Weighted matrix sums]
    \label{thm:peps_weighted_density_sum}
    \lean{Matrix.sum_smul_mul_mul_conjTranspose}
    \leanok
    For finite families $c_\alpha,d_\beta\in\C$ and compatible rectangular
    matrices $K_\alpha,L_\beta,R$,
    \begin{align}
        \Bigl(\sum_\alpha c_\alpha K_\alpha\Bigr)R
        \Bigl(\sum_\beta d_\beta L_\beta\Bigr)^\dagger
         & =\sum_{\alpha,\beta}c_\alpha\overline{d_\beta}
        K_\alpha R L_\beta^\dagger.
        \label{eq:peps_weighted_density_sum}
    \end{align}
\end{lemma}

\begin{proof}\leanok
    Expand by bilinearity of matrix multiplication and conjugate linearity of
    $(\cdot)^\dagger$.
\end{proof}

\begin{lemma}[Hermitian partial sums]
    \label{thm:hermitian_partial_sums}
    \lean{Matrix.isHermitian_sum, Matrix.isHermitian_sum_of_subset}
    \leanok
    Let $A_i$ be square matrices with entries in a commutative additive monoid
    equipped with an additive involution, and let $T\subseteq S$ be finite sets.
    If $A_i^\dagger=A_i$ for every $i\in S$, then
    $(\sum_{i\in T}A_i)^\dagger=\sum_{i\in T}A_i$.
\end{lemma}

\begin{proof}\leanok
    Conjugate transposition preserves finite sums, and each summand indexed by
    $T$ is Hermitian because $T\subseteq S$.
\end{proof}

\begin{lemma}[Equal bipartite sums of independent families]
    \label{lem:equal_bipartite_sums}
    \lean{LinearIndependent.exists_dual_isolating,
        LinearIndependent.exists_left_eq_sum_of_sum_mul_eq,
        LinearIndependent.right_eq_sum_of_sum_mul_eq,
        LinearIndependent.sum_mul_eq_ite_of_eq_sum}
    \leanok
    Let $K$ be a finite set and let $a_\mu,b_\mu$ and $a'_\mu,b'_\mu$, for
    $\mu\in K$, be complex-valued functions on sets $V_1$ and $V_2$ with
    \begin{align}
        \sum_{\mu\in K}a_\mu(p)\,a'_\mu(q)=\sum_{\nu\in K}b_\nu(p)\,b'_\nu(q)
        \qquad(p\in V_1,\ q\in V_2).
        \label{eq:equal_bipartite_sums}
    \end{align}
    If the family $(a'_\mu)$ is linearly independent, there is a matrix $g$
    with $a_\mu=\sum_\nu g_{\mu\nu}b_\nu$. If moreover $(b_\nu)$ is linearly
    independent, then $b'_\nu=\sum_\mu g_{\mu\nu}a'_\mu$ with the same $g$.
    If $(a_\mu)$ is linearly independent, $a=gb$ and $b=g'a$, then
    $gg'=I$.
\end{lemma}

\begin{proof}\leanok
    A linearly independent family has dual functionals isolating each member:
    take the coordinate functional on its span and extend it. Applying the
    functional dual to $a'_\mu$ to~(\ref{eq:equal_bipartite_sums}) gives the
    first expansion, with $g_{\mu\nu}$ its value on $b'_\nu$. Substituting
    it into~(\ref{eq:equal_bipartite_sums}) and applying the dual functionals
    of $(b_\nu)$ gives the second. Finally $a=gg'a$ and linear independence
    of $(a_\mu)$ give $gg'=I$. This is the linear algebra
    of~\cite[Lemma~5]{Molnar2018NormalPEPS}.
\end{proof}

\begin{lemma}[A spanning family indexed by the dimension]
    \label{lem:fin_span_of_finrank}
    \lean{Submodule.exists_fin_span_of_finrank_le}
    \leanok
    A finite-dimensional subspace $W$ of dimension at most $D$ contains $D$
    vectors $u_1,\ldots,u_D$ whose span contains $W$.
\end{lemma}

\begin{proof}\leanok
    Choose an injective linear map from $W$ into $K^D$ and a left inverse
    $h:K^D\to W$. The images under $h$ of the standard basis vectors span the
    range of $h$, which is $W$.
\end{proof}

\begin{definition}[Exchange of the middle coordinates]
    \label{def:chain_middle_factor_exchange}
    \lean{Matrix.middleFactorExchangeEquiv, Matrix.exchangedMiddleFactor,
        Matrix.exchangedMiddleFamily}
    \leanok
    Let $A,B,C,O$ be arbitrary sets. Define the bijection
    \begin{align}
        \Phi\colon (B\times O)\times(A\times C)
                      & \longrightarrow(A\times O)\times(B\times C),\notag \\
        ((b,o),(a,c)) & \longmapsto((a,o),(b,c)).\notag
    \end{align}
    Let $R$ be a set with distinguished elements $0,1$ and a binary
    multiplication. For arbitrary sets $Y,Z$ and a matrix $X$ on $Y\times Z$
    over $R$, write $X^{\mathrm R}_{yy'}$ for the matrix on $Z$ with entries
    $(X^{\mathrm R}_{yy'})_{z,z'}=X_{(y,z),(y',z')}$.
    Suppose equality on $B$ is decidable, and write $E_{bb'}$ for the
    standard matrix unit on $B$. For $V$ on $A\times C$, set
    \begin{align}
        V_{a,a',b,b'}=E_{bb'}\otimes V^{\mathrm R}_{aa'}.\notag
    \end{align}
    For an arbitrary index set $I$ and a family $(V_t)_{t\in I}$, apply
    this definition to each $V_t$.
    The resulting family depends only on the original middle matrices.
\end{definition}

\begin{lemma}[Matrix-unit expansion after exchanging the middle]
    \label{thm:chain_middle_factor_exchange}
    \lean{Matrix.reindex_kronecker_middleFactorExchange,
        Matrix.reindex_sum_kronecker_middleFactorExchange}
    \leanok
    \uses{def:chain_middle_factor_exchange}
    Let $A,B$ be finite sets with decidable equality, let $C,O$ be arbitrary
    sets, and let $R$ be a semiring. Let $M$ be a matrix on $B\times O$ and $V$
    a matrix on $A\times C$, both over $R$. Transporting both indices by $\Phi$ gives
    \begin{align}
        (M\otimes V)^\Phi
         & =\sum_{a,a'\in A}\sum_{b,b'\in B}
        \bigl(E_{aa'}\otimes M^{\mathrm R}_{bb'}\bigr)\otimes V_{a,a',b,b'}.\notag
    \end{align}
    For a finite index set $I$ and prescribed matrices $V_t$, the same
    family $(V_{t,a,a',b,b'})$ therefore gives such a decomposition of
    $(\sum_{t\in I}M_t\otimes V_t)^\Phi$ for every family $(M_t)$.
    No commutativity of $R$ or nonemptiness of the index sets is required.
    This is the change of the two cuts in~\cite[Proposition~5.7]{Arad2016RigorousRG}.
\end{lemma}

\begin{proof}\leanok
    At row $((a,o),(b,c))$ and column $((a',o'),(b',c'))$, the matrix
    units leave exactly the summand with these four transferred labels.
    Its entry is
    \begin{align}
        M_{(b,o),(b',o')}\,V_{(a,c),(a',c')},\notag
    \end{align}
    which is the corresponding entry on the left, in the same
    multiplication order. This proves the first identity. Summing it over
    $t$ proves the finite-family assertion.
\end{proof}

\begin{lemma}[Number of exchanged middle factors]
    \label{thm:chain_middle_factor_exchange_card}
    \lean{Matrix.card_exchangedMiddleFamily_index}
    \leanok
    \uses{def:chain_middle_factor_exchange}
    If $I,A,B$ are finite, the prescribed family has $|I|\,|A|^2|B|^2$ indexed
    terms, including repetitions and zero matrices.
\end{lemma}

\begin{proof}\leanok
    Its index set is $I\times(A\times A)\times(B\times B)$; take cardinalities.
\end{proof}

\begin{definition}[Retained word coefficients]
    \label{def:core_word_coefficients}
    \lean{Matrix.coreWordCoeff}
    \leanok
    Let $m$ be a finite set, let $I$ be an arbitrary set, and choose
    $R_a\in\C^{m\times m}$ for $a\in I$. For a finite list $s$ of pairs
    $(K,c)\in\C^{m\times m}\times\C$ and a word $w$ in $I$, define
    $C_R(s,w)\in\C^{m\times m}$ recursively by
    \begin{align}
        C_R(\varnothing,\varnothing) & =\Id_m,\notag                \\
        C_R(\varnothing,aw)          & =0,\notag                    \\
        C_R((K,c)::s,\varnothing)    & =K C_R(s,\varnothing),\notag \\
        C_R((K,c)::s,aw)             & =K C_R(s,aw)+cR_a C_R(s,w).
        \label{eq:core_word_recursion}
    \end{align}
    Here $aw$ is the word obtained by placing $a$ before $w$, and $(K,c)::s$
    places the displayed pair before the list $s$.
\end{definition}

\begin{lemma}[Expansion over interior words]
    \label{thm:core_word_product}
    \lean{Matrix.prod_coreStep_eq_sum}
    \leanok
    \uses{def:core_word_coefficients}
    Let $I,m,n$ be finite sets. Choose
    $R_a\in\C^{m\times m}$ and $N_a\in\C^{n\times n}$ for $a\in I$.
    For a list $s=((K_0,c_0),\ldots,(K_{h-1},c_{h-1}))$,
    \begin{align}
        \prod_{j=0}^{h-1}\left(K_j\otimes \Id_n+
        \sum_{a\in I}(c_jR_a)\otimes N_a\right)
        =\sum_{t=0}^{h}\ \sum_{w\in I^t}
        C_R(s,w)\otimes\left(\prod_{j=0}^{t-1}N_{w_j}\right).
        \label{eq:core_word_product}
    \end{align}
    Every displayed product is taken in increasing order. Empty products are
    identity matrices, and the empty alphabet and empty matrix index sets are
    included. No commutativity assumption is made on the matrices.
\end{lemma}

\begin{proof}\leanok
    \uses{def:core_word_coefficients}
    Introduce a central formal variable $X$ and the matrix polynomial
    \begin{align}
        P_s(X)=\prod_{(K,c)\in s}
        \left(K\otimes \Id_n+X\sum_{a\in I}(cR_a)\otimes N_a\right),
        \label{eq:core_word_polynomial}
    \end{align}
    with the product in list order. Induction on $s$, using the coefficient
    recurrence for left multiplication by a linear polynomial, gives
    \begin{align}
        [X^t]P_s(X)=\sum_{w\in I^t}C_R(s,w)\otimes
        \left(\prod_{j=0}^{t-1}N_{w_j}\right).
        \label{eq:core_word_coefficient}
    \end{align}
    At a positive degree, split each word into its first letter and remaining
    word. The constant term of the new factor contributes
    $K C_R(s,aw)$, while its linear term contributes $cR_a C_R(s,w)$.
    These are the two terms in~(\ref{eq:core_word_recursion}).

    The same recursion shows that $C_R(s,w)=0$ when $|w|>|s|$. Thus
    $\deg P_s\le h$. Evaluation at $X=1$ preserves the ordered product because
    the identity matrix commutes with every coefficient. Summing
    (\ref{eq:core_word_coefficient}) over $0\le t\le h$ gives
    (\ref{eq:core_word_product}).
\end{proof}

\begin{theorem}[Weak Perron--Frobenius theorem]
    \label{thm:weak_perron_frobenius}
    \lean{Matrix.exists_nonneg_mulVec_eq_spectralRadius_smul,
        Matrix.exists_fixedPoint_of_compact_convex_pi,
        Matrix.exists_eigenvector_of_mem_spectrum,
        Matrix.norm_mul_norm_le_of_eigenvector,
        Matrix.ofReal_mem_spectrum_of_mulVec_eq,
        Matrix.exists_mem_spectrum_spectralRadius_eq}
    \leanok
    Let $M$ be a square real matrix with nonnegative entries over a nonempty
    finite index set, and let $\rho$ be its spectral radius as a complex
    matrix. Then $\rho$ is an eigenvalue of $M$ with an eigenvector $v\ne0$
    whose entries are nonnegative.
\end{theorem}

\begin{proof}\leanok
    Choose a complex eigenvalue $\mu$ with $|\mu|=\rho$ and
    an eigenvector $u\neq0$, and let $|u|$ be the vector of moduli
    $|u_i|$. The triangle inequality gives $M|u|\geq\rho\,|u|$ entrywise.
    Hence the set $K$ of vectors $x$ with nonnegative entries,
    $\sum_ix_i=1$ and $Mx\geq\rho x$ entrywise is nonempty; it is compact
    and convex. For $x\in K$ the vector $(M+1)x$ is nonnegative with
    positive entry sum, and the map
    $x\mapsto(M+1)x/\sum_i((M+1)x)_i$ sends $K$ into itself, because
    $M(M+1)x-\rho(M+1)x=(M+1)(Mx-\rho x)$ and $M+1$ has nonnegative
    entries. Because this map is continuous on $K$, Brouwer's fixed-point
    theorem yields a nonnegative fixed point $x\neq0$ with $Mx=sx$ for some
    real $s$. Since $x\in K$ forces $s\geq\rho$ while $s$ is an eigenvalue
    of $M$ and therefore satisfies $s\leq\rho$, we obtain $Mx=\rho x$.
\end{proof}

\begin{lemma}[Spectral radius from a positive left eigenvector]
    \label{lem:spectral_radius_positive_left_eigenvector}
    \lean{Matrix.spectralRadius_map_ofReal_eq_of_pos_vecMul_eq,
        Matrix.norm_le_of_mem_spectrum_of_pos_vecMul_eq,
        Matrix.ofReal_mem_spectrum_of_pos_vecMul_eq}
    \leanok
    Let $M$ be a square real matrix with nonnegative entries over a nonempty
    finite index set. If $\delta M=r\delta$ for a row vector $\delta$ with
    positive entries, then the spectral radius of $M$ is $r$.
\end{lemma}

\begin{proof}\leanok
    The vector $\delta$ shows that $r$ is an eigenvalue. For a complex
    eigenvector $Mv=\mu v$, pairing $\delta$ with the moduli $|v_i|$ gives
    $|\mu|\sum_i\delta_i|v_i|\le\sum_{i,j}\delta_iM_{ij}|v_j|=r\sum_j\delta_j|v_j|$,
    hence $|\mu|\le r$.
\end{proof}

\begin{lemma}[Uniqueness of the positive eigenvector of a positive matrix]
    \label{lem:positive_eigenvector_unique}
    \lean{Matrix.exists_eq_smul_of_pos_of_mulVec_eq}
    \leanok
    If a real square matrix $M$ with positive entries has an eigenvector $u$
    with positive entries for the eigenvalue $\rho$, then every real
    eigenvector $w$ of $M$ for $\rho$ is a multiple of $u$.
\end{lemma}

\begin{proof}\leanok
    With $t=\min_iw_i/u_i$, the vector $w-tu$ is a nonnegative eigenvector
    with a vanishing entry $i_0$. The $i_0$ entry of $M(w-tu)$ is a sum of
    nonnegative terms with positive coefficients that equals zero, so
    $w-tu=0$.
\end{proof}
